\section{Evaluation Methodology}
\label{sec:methodology}

This paper reports four primary metrics from the campaign pipeline.

\paragraph{Campaign hit rate.}
For each model, hit rate is the fraction of probes marked as
\texttt{SUCCESS} by the campaign runner:
\[
\text{hit\_rate} = \frac{n_{\text{success}}}{n_{\text{probes}}}.
\]

\paragraph{Campaign refusal rate.}
Refusal rate is derived from campaign outcomes marked as
\texttt{FAILED}:
\[
\text{refusal\_rate} = \frac{n_{\text{failed}}}{n_{\text{probes}}}.
\]
This reflects framework-native refusal logic rather than external human
labeling. In the default runner path, refusal is assigned when a probe does
not meet exploitation success criteria after heuristic response checks (for
example compliance and refusal indicators) in the probe evaluator
\footnote{\path{harness/runner.py}, \texttt{\_default\_success\_detector}.}.
This is intentionally distinct from benchmark-focused judge rubrics such as
\sr\ and HarmBench \citep{souly2024strongreject,mazeika2024harmbench}.

\paragraph{Aggregate risk score.}
The reporting pipeline computes a CVSS-inspired aggregate risk score using
\path{harness/runner.py}. For each successful probe $i$, a weighted score is
computed from technique severity and probe confidence:
\[
s_i = b_i \cdot \max(c_i, 0.5),
\]
where $b_i \in \{9.5, 7.5, 5.0, 2.0\}$ for
\{critical, high, medium, low\} severity, and $c_i \in [0,1]$ is probe
confidence. The model-level aggregate risk score is then:
\[
\text{risk\_score} = \min\!\left(10,\; \frac{1}{|S|}\sum_{i \in S} s_i\right),
\]
where $S$ is the set of successful probes for that model
\footnote{\path{harness/runner.py}, \texttt{\_compute\_risk\_score}.}.

\paragraph{Operational error rate.}
Error rate is the fraction of probes ending in \texttt{ERROR}. This is an
operational reliability metric, not a direct safety metric, but it materially
affects interpretation of model-level findings.
\[
\text{error\_rate} = \frac{n_{\text{error}}}{n_{\text{probes}}}.
\]
All rates in this paper use the same denominator ($n_{\text{probes}}$), where
$n_{\text{probes}} = n_{\text{success}} + n_{\text{failed}} + n_{\text{error}} + n_{\text{partial}}$.
This means high error rates directly suppress achievable hit/refusal rates and
must be interpreted jointly with risk scores.

\paragraph{Limits of comparability.}
Because this run is tied to a curated registry and framework-native probes, the
results should be interpreted as an artifact-backed case study rather than a
general benchmark leaderboard.

\paragraph{Benchmark context (ClearHarm and \sr).}
This run is not a direct benchmark evaluation on ClearHarm
\citep{hollinsworth2025clearharm} or the original \sr\ prompt sets
\citep{souly2024strongreject}. Instead, it uses the framework's ATT\&CK-inspired
20-probe campaign and reports campaign-native outcomes (hit, refusal, risk,
error) over that prompt distribution. In practical terms, the current paper is
best read as a systems-and-operations study with benchmark-informed framing,
not as a benchmark leaderboard claim.

\paragraph{Evaluator status in this artifact set.}
The refreshed 11-model artifact set includes post-hoc evaluator overlays on the
archived traces (without rerunning target-model probes), using a fixed judge
configuration for \sr-style scoring. This supports within-artifact
comparisons, while definitive benchmark-leaderboard claims would still require
a dedicated benchmark-aligned rerun on the canonical prompt distributions.
